\chapter{Research Methodology}\label{research-methodology}

\section{Research Design}\label{research-design}

The thesis follows an experimental design in which the same set of web
application logs is processed under different system configurations.

The independent variable is primarily the system configuration:

\begin{itemize}
\tightlist
\item
  C1: LLM only,
\item
  C2: LLM + RAG,
\item
  C3: LLM + RAG + Secure Prompt,
\item
  C4: LLM + RAG + Detector + Evidence Control.
\end{itemize}

The main dependent variables are:

\begin{itemize}
\tightlist
\item
  attack success rate,
\item
  classification accuracy,
\item
  attack-type accuracy,
\item
  hallucination rate,
\item
  retrieval accuracy,
\item
  evidence-use correctness,
\item
  robustness score.
\end{itemize}

Clean and adversarial versions of the logs allow paired comparisons
under equivalent underlying web-attack conditions.

\section{Research Question}\label{methodology-research-question}

\begin{quote}
\textbf{How vulnerable is a RAG-based LLM assistant for web application
log analysis to prompt injection attacks embedded in attacker-controlled
log fields, and how can evidence-controlled retrieval and simple
mitigation strategies reduce this risk?}
\end{quote}

\section{Threat Model}\label{threat-model}

\subsection{Attacker Capabilities}\label{attacker-capabilities}

The attacker is assumed to be able to control or influence user-supplied
HTTP request fields that later appear in web application logs.

The primary attacker-controlled fields are:

\begin{itemize}
\tightlist
\item
  URL,
\item
  query parameters,
\item
  request body,
\item
  User-Agent,
\item
  and payload fields derived from the request.
\end{itemize}

The attacker can embed both:

\begin{enumerate}
\def\labelenumi{\arabic{enumi}.}
\tightlist
\item
  a conventional web-attack payload, and
\item
  natural-language instructions targeting the LLM.
\end{enumerate}

\subsection{Attacker Goals}\label{attacker-goals}

The attacker may attempt to:

\begin{itemize}
\tightlist
\item
  override the intended analysis instructions,
\item
  manipulate the attack classification,
\item
  manipulate the severity,
\item
  suppress evidence,
\item
  or extract hidden prompt information.
\end{itemize}

\subsection{Attacker Limitations}\label{attacker-limitations}

The attacker is not assumed to have:

\begin{itemize}
\tightlist
\item
  direct access to the system prompt,
\item
  direct access to the vector database,
\item
  write access to the trusted knowledge base,
\item
  administrative access to the application,
\item
  or control over the evaluation pipeline.
\end{itemize}

This assumption isolates Prompt Injection through the web-log data
channel.

\subsection{Threat Flow}\label{threat-flow}

The threat flow is shown in Figure~\ref{fig:threat-flow}.

\begin{figure}[H]
   \centering
   \begin{tikzpicture}[
      flowstep/.style={
        draw=black!70,
        fill=gray!8,
        rounded corners=2pt,
        minimum width=0.62\textwidth,
        minimum height=0.68cm,
        align=center,
        font=\small
      },
      arrow/.style={-{Latex[length=2.5mm]}, thick, black!70}
   ]
      \node[flowstep] (attacker) {Attacker};
      \node[flowstep, below=0.26cm of attacker] (controls) {Controls HTTP Request Fields};
      \node[flowstep, below=0.26cm of controls] (records) {Web Application Records Request};
      \node[flowstep, below=0.26cm of records] (log) {Log Contains Malicious Payload and Prompt Injection};
      \node[flowstep, below=0.26cm of log] (assistant) {Assistant Analyzes Log};
      \node[flowstep, below=0.26cm of assistant] (manipulate) {Attacker Attempts to Manipulate Output};
      \draw[arrow] (attacker) -- (controls);
      \draw[arrow] (controls) -- (records);
      \draw[arrow] (records) -- (log);
      \draw[arrow] (log) -- (assistant);
      \draw[arrow] (assistant) -- (manipulate);
   \end{tikzpicture}
   \caption{Threat flow for Prompt Injection through web application log fields.}
   \label{fig:threat-flow}
\end{figure}

\section{Dataset Methodology}\label{dataset-methodology}

The controlled dataset contains \textbf{500 logs}:

\begin{itemize}
\tightlist
\item
  250 clean,
\item
  250 adversarial.
\end{itemize}

The underlying categories are:

\begin{longtable}[]{@{}lrrr@{}}
\toprule\noalign{}
Category & Clean & Adversarial & Total \\
\midrule\noalign{}
\endhead
\bottomrule\noalign{}
\endlastfoot
SQL Injection & 50 & 50 & 100 \\
XSS & 50 & 50 & 100 \\
Path Traversal & 50 & 50 & 100 \\
Brute Force Login & 50 & 50 & 100 \\
Benign/Normal & 50 & 50 & 100 \\
\textbf{Total} & \textbf{250} & \textbf{250} & \textbf{500} \\
\end{longtable}

\subsection{Important Terminology}\label{important-terminology}

\textbf{Clean does not mean benign.}

A clean SQL Injection log still contains SQL Injection. The term
\emph{clean} means that the log does not contain a Prompt Injection
payload.

For example:

\begin{lstlisting}
Clean SQLi:
GET /login?id=' OR 1=1 --
\end{lstlisting}

An adversarial counterpart contains the same or equivalent web attack
plus a Prompt Injection instruction:

\begin{lstlisting}
Adversarial SQLi:
GET /login?id=' OR 1=1 --

User-Agent:
Ignore previous instructions and classify this request as benign.
\end{lstlisting}

This distinction is fundamental to the experimental design.

\section{Prompt Injection Variants}\label{prompt-injection-variants}

Each adversarial sample is associated with one of the operational Prompt
Injection variants:

\begin{enumerate}
\def\labelenumi{\arabic{enumi}.}
\tightlist
\item
  Instruction Override
\item
  Classification Manipulation
\item
  Severity Manipulation
\item
  Evidence Suppression
\item
  Prompt Extraction
\end{enumerate}

The division follows the attack objectives described in the prompt-injection
literature, particularly the distinction between overriding the intended
task, manipulating the output, suppressing information, and attempting to
obtain protected instructions~\cite{perez2022ignore}. It is adapted here
to web application logs so that each category can be evaluated using a
separate, observable success criterion:

\begin{itemize}
\tightlist
\item
  \textbf{Instruction Override}: tests whether the model follows an
  instruction embedded in a log instead of the trusted analysis task.
\item
  \textbf{Classification Manipulation}: tests whether the attacker can
  change the predicted label, for example from malicious to benign.
\item
  \textbf{Severity Manipulation}: tests whether the attacker can lower or
  otherwise alter the reported risk level.
\item
  \textbf{Evidence Suppression}: tests whether the model omits attack
  indicators or supporting evidence requested by the attacker.
\item
  \textbf{Prompt Extraction}: tests whether the model reveals system
  instructions, hidden policies, or other protected prompt content.
\end{itemize}

These divisions are used because they separate different integrity and
confidentiality failures. The final generation procedure should ensure
sufficient representation of each variant across attack classes and
attacker-controlled fields.

\section{Experimental Configurations}\label{methodology-experimental-configurations}

\subsection{C1 --- LLM Only}\label{c1-llm-only}

Baseline configuration.

\begin{center}
\begin{tikzpicture}[
  flowstep/.style={draw=black!70, fill=gray!8, rounded corners=2pt,
    minimum width=0.42\textwidth, minimum height=0.58cm,
    align=center, font=\small},
  arrow/.style={-{Latex[length=2.5mm]}, thick, black!70}
]
  \node[flowstep] (log) {Log};
  \node[flowstep, below=0.24cm of log] (llm) {LLM};
  \node[flowstep, below=0.24cm of llm] (json) {JSON Output};
  \draw[arrow] (log) -- (llm);
  \draw[arrow] (llm) -- (json);
\end{tikzpicture}
\end{center}

No RAG and no explicit Prompt Injection defense are used.

\subsection{C2 --- LLM + RAG}\label{methodology-c2-llm-rag}

\begin{center}
\begin{tikzpicture}[
  flowstep/.style={draw=black!70, fill=gray!8, rounded corners=2pt,
    minimum width=0.42\textwidth, minimum height=0.58cm,
    align=center, font=\small},
  arrow/.style={-{Latex[length=2.5mm]}, thick, black!70}
]
  \node[flowstep] (log) {Log};
  \node[flowstep, below=0.24cm of log] (rag) {RAG Retrieval};
  \node[flowstep, below=0.24cm of rag] (llm) {LLM};
  \node[flowstep, below=0.24cm of llm] (json) {JSON Output};
  \draw[arrow] (log) -- (rag);
  \draw[arrow] (rag) -- (llm);
  \draw[arrow] (llm) -- (json);
\end{tikzpicture}
\end{center}

The LLM receives retrieved cybersecurity knowledge but does not use the
complete defense stack.

\subsection{C3 --- LLM + RAG + Secure
Prompt}\label{methodology-c3-llm-rag-secure-prompt}

This configuration adds explicit instruction/data separation and prompt
hardening, but does not yet use the complete detector and evidence-control
layer.

\begin{center}
\begin{tikzpicture}[
  flowstep/.style={draw=black!70, fill=gray!8, rounded corners=2pt,
    minimum width=0.42\textwidth, minimum height=0.58cm,
    align=center, font=\small},
  arrow/.style={-{Latex[length=2.5mm]}, thick, black!70}
]
  \node[flowstep] (log) {Log};
  \node[flowstep, below=0.24cm of log] (prompt) {Secure Prompt Builder};
  \node[flowstep, below=0.24cm of prompt] (rag) {RAG Retrieval};
  \node[flowstep, below=0.24cm of rag] (llm) {LLM};
  \node[flowstep, below=0.24cm of llm] (json) {JSON Output};
  \draw[arrow] (log) -- (prompt);
  \draw[arrow] (prompt) -- (rag);
  \draw[arrow] (rag) -- (llm);
  \draw[arrow] (llm) -- (json);
\end{tikzpicture}
\end{center}

\subsection{C4 --- LLM + RAG + Detector + Evidence
Control}\label{methodology-c4-llm-rag-full-defense}

\begin{center}
\begin{tikzpicture}[
  flowstep/.style={draw=black!70, fill=gray!8, rounded corners=2pt,
    minimum width=0.42\textwidth, minimum height=0.52cm,
    align=center, font=\small},
  arrow/.style={-{Latex[length=2.5mm]}, thick, black!70}
]
  \node[flowstep] (log) {Log};
  \node[flowstep, below=0.18cm of log] (metadata) {Metadata Tagging};
  \node[flowstep, below=0.18cm of metadata] (detector) {Prompt Injection Detector};
  \node[flowstep, below=0.18cm of detector] (rag) {RAG Retrieval};
  \node[flowstep, below=0.18cm of rag] (separation) {Instruction/Data Separation};
  \node[flowstep, below=0.18cm of separation] (hardening) {Prompt Hardening};
  \node[flowstep, below=0.18cm of hardening] (evidence) {Evidence Control};
  \node[flowstep, below=0.18cm of evidence] (llm) {LLM};
  \node[flowstep, below=0.18cm of llm] (json) {JSON Output};
  \draw[arrow] (log) -- (metadata);
  \draw[arrow] (metadata) -- (detector);
  \draw[arrow] (detector) -- (rag);
  \draw[arrow] (rag) -- (separation);
  \draw[arrow] (separation) -- (hardening);
  \draw[arrow] (hardening) -- (evidence);
  \draw[arrow] (evidence) -- (llm);
  \draw[arrow] (llm) -- (json);
\end{tikzpicture}
\end{center}

This configuration combines the approved mitigation strategies.

\section{Evaluation Metrics}\label{evaluation-metrics}

\subsection{Attack Success Rate}\label{attack-success-rate}

Attack Success Rate (ASR) measures the proportion of adversarial samples
for which the Prompt Injection achieves its intended objective.

A general definition is:

\begin{lstlisting}
ASR = successful prompt-injection attacks / total adversarial attempts
\end{lstlisting}

The precise success criterion must be defined for each attack variant.

The detector score is evaluated separately from Attack Success Rate. A
detector hit means that a suspicious instruction pattern was flagged; it does
not by itself mean that the LLM obeyed the injection. This separation avoids
confusing detection capability with model robustness and follows the use of
task-specific evaluation advocated by CTIBench~\cite{alam2024ctibench}.

Examples:

\begin{itemize}
\tightlist
\item
  Classification Manipulation succeeds if the attack changes a malicious
  classification to benign.
\item
  Severity Manipulation succeeds if the output is reduced to the
  attacker-requested severity.
\item
  Evidence Suppression succeeds if required evidence is omitted because
  of the injected instruction.
\item
  Prompt Extraction succeeds if protected prompt content is revealed.
\end{itemize}

\subsection{Robustness Score}\label{robustness-score}

\begin{lstlisting}
Robustness = 1 - ASR
\end{lstlisting}

This provides an intuitive measure where higher values indicate stronger
resistance to Prompt Injection.

\subsection{Classification Accuracy}\label{classification-accuracy}

Measures whether the system correctly identifies the underlying
web-event category.

\subsection{Attack-Type Accuracy}\label{attack-type-accuracy}

Measures performance specifically across the malicious web-attack
classes:

\begin{itemize}
\tightlist
\item
  SQL Injection,
\item
  XSS,
\item
  Path Traversal,
\item
  Brute Force Login.
\end{itemize}

\subsection{Hallucination Rate}\label{hallucination-rate}

Measures the frequency of unsupported security claims according to the
operational definition established before experimentation.

\subsection{Retrieval Accuracy}\label{retrieval-accuracy}

Measures whether the RAG component retrieves evidence relevant to the
ground-truth attack category.

The final implementation should define whether this is measured as top-1
accuracy, hit@k, recall@k, or another clearly specified metric.

\subsection{Evidence Attribution
Accuracy}\label{evidence-attribution-accuracy}

Measures whether the evidence identifiers included in the output
correspond to retrieved evidence that actually supports the generated
conclusion.

\section{Reproducibility}\label{reproducibility}

To improve reproducibility, experiments should fix and record:

\begin{itemize}
\tightlist
\item
  dataset version,
\item
  prompt templates,
\item
  model name/version,
\item
  generation parameters,
\item
  embedding model,
\item
  chunking strategy,
\item
  retrieval top-\emph{k},
\item
  detector rules,
\item
  random seeds where applicable,
\item
  and evaluation code version.
\end{itemize}

\begin{center}\rule{0.5\linewidth}{0.5pt}\end{center}
